%!TEX root = ../dolgozat.tex
\chapter{Összefoglalás}
\label{ch:osszefoglalas}

Dolgozatomban egy olyan webes vizsgáztató rendszer prototípusát terveztem meg és készítettem el,
amely egy vizsgán belül vegyes feladattípusokat kezel, a programozási feladatokat elszigetelt
Docker-konténerben, tesztesetekkel automatikusan értékeli, és valós időben jelzi az oktatónak, ha
egy hallgató a vizsga közben elhagyja a vizsgaablakot.

A szakirodalom áttekintése során megismertem a számítógéppel támogatott értékelés lehetőségeit
és korlátait, a programozási feladatok automatikus értékelésének több évtizedes múltját, valamint
az online vizsgafelügyelet körüli szakmai és etikai vitát. A hasonló rendszerek vizsgálata azt
mutatta, hogy a vegyes feladattípusok, a biztonságos kódfuttatás és a valós idejű felügyelet
ritkán jelenik meg együtt, és ahol igen, ott jelentős üzemeltetési terhet jelent. Ezekre a
tapasztalatokra építve fogalmaztam meg a rendszer követelményeit.

A megvalósított rendszer hat feladattípust támogat. A zárt és a rövid szöveges feladatokat
automatikusan pontozza, a több helyes választ tartalmazó feladatoknál a hibás jelöléseket levonó
részpontozással, a rövid válaszoknál pedig a formai eltéréseket (kis- és nagybetűk, szóközök,
tizedesjel) figyelmen kívül hagyó összehasonlítással. A programozási feladatokat nyilvános és
rejtett tesztesetekkel értékeli, a kifejtős feladatokhoz pedig kézi javítási felületet és
visszajelzési lehetőséget biztosít. A vizsgafolyamatban a határidőt a szerver érvényesíti, a
válaszok automatikusan mentődnek, és a hallgató munkája akkor sem vész el, ha a vizsga közben
bezárja a böngészőt.

A tesztelés eredményei alapján a rendszer a kitűzött céloknak megfelel. A 32 funkcionális
teszteset mindegyike, valamint a valós idejű felügyelet 10 tesztesete is sikeres volt. A sandbox
biztonsági tesztjei igazolták, hogy a tervezett korlátozások a gyakorlatban is megállítják a
leggyakoribb támadásokat: a végtelen ciklust, a fork bombát, a túlzott memóriafoglalást, a hálózati
és a fájlrendszer-hozzáférést. Egy programozási feladat értékelése 20 teszteset esetén is
nagyjából egy másodperc alatt elkészül, ami annak köszönhető, hogy egy kiértékeléshez csak
egyetlen konténer indul. A felügyeleti jelzés helyi környezetben 10 milliszekundumon belül
eljutott az oktatóhoz.

Amit kevésbé sikerültnek tartok, az elsősorban a kódfuttatás eredményének feldolgozása: a
biztonsági tesztek során kiderült, hogy egy, az eredményjelölő szövegét kiíró program
megakaszthatta a saját értékelését. Bár ez a pontszám meghamisítására nem adott lehetőséget, és a
hibát véletlenszerű, futtatásonként generált jelölővel ki is javítottam, jól mutatja, hogy egy nem megbízható kódot futtató rendszernél a vezérlő és a
futtatott program közötti kommunikációt is ugyanolyan gondosan kell megtervezni, mint magát az
elszigetelést. Ez a tapasztalat egyben megerősítette bennem, hogy egy ilyen rendszernél a
szándékosan ártó bemenetekkel végzett tesztelés elengedhetetlen. A felügyelet terén pedig
tudatosan választottam a korlátozott, de a magánszférát tiszteletben tartó megoldást; véleményem
szerint egy egyetemi vizsgán ez jobb egyensúlyt jelent, mint a hallgatók kamerás megfigyelése,
ugyanakkor el kell fogadni, hogy egy elszánt hallgatót ez a módszer nem tart vissza.

A rendszer több irányban is továbbfejleszthető. A legkézenfekvőbb lépés további programozási
nyelvek (például C, C++ vagy Java) támogatása, amely a sandbox felépítésének köszönhetően
elsősorban a megfelelő konténerképek és a fordítási lépés hozzáadását jelentené. A kód
szerkesztését egy szintaxiskiemelést és automatikus kiegészítést nyújtó szerkesztő, például a
Monaco Editor beépítése tenné kényelmesebbé. Nagyobb létszámú vizsgák esetén a kódértékelést
egy üzenetsorra és külön feldolgozó folyamatokra (\textit{worker}) érdemes áthelyezni, hogy a
kiértékelés ne terhelje a webszervert, és a terhelés több gépre is szétosztható legyen. A
vizsgák tisztaságát a beküldött programok egymás közötti hasonlóságát vizsgáló plágiumkeresés
segíthetné, az oktatók munkáját pedig a feladatonkénti statisztikák (átlagpontszám, a sikertelen
tesztesetek eloszlása) és egy feladatbank, amelyből a feladatok több vizsgában is
újrahasznosíthatók. Végül egy LTI szabvány szerinti integráció lehetővé tenné, hogy a rendszer
egy meglévő tanulásirányítási rendszerbe, például a Moodle-be illeszkedjen, így a hallgatóknak nem
kellene külön fiókot létrehozniuk.

Összességében úgy érzem, hogy a dolgozat elkészítése során nemcsak egy működő rendszert hoztam
létre, hanem a webes alkalmazásfejlesztés, az informatikai biztonság és a konténerizáció terén
is olyan gyakorlati tapasztalatokat szereztem, amelyeket a későbbi szakmai munkám során is
hasznosítani tudok.
